\section*{Capítulo \ref{ch:b3:complex-spectra-magnetism} --- Espectros electrónicos y magnetismo de los complejos}

\begin{solution}{exo:b3:complex-spectra-magnetism:1}
D: E$_g$ + T$_{2g}$ ($2 + 3 = 5$); F: A$_{2g}$ + T$_{1g}$ + T$_{2g}$ ($1 + 3 + 3 = 7$); G: A$_{1g}$ + E$_g$ + T$_{1g}$ + T$_{2g}$
($1 + 2 + 3 + 3 = 9$).
\end{solution}

\begin{solution}{exo:b3:complex-spectra-magnetism:2}
$d^1$ $^2$D $\to$ $^2$T$_{2g}$; $d^2$ $^3$F $\to$ $^3$T$_{1g}$; $d^3$ $^4$F $\to$ $^4$A$_{2g}$; $d^4$ $^5$D $\to$ $^5$E$_g$; $d^5$ $^6$S $\to$ $^6$A$_{1g}$; $d^6$ $^5$D
$\to$ $^5$T$_{2g}$; $d^7$ $^4$F $\to$ $^4$T$_{1g}$; $d^8$ $^3$F $\to$ $^3$A$_{2g}$; $d^9$ $^2$D $\to$ $^2$E$_g$.
\end{solution}

\begin{solution}{exo:b3:complex-spectra-magnetism:3}
\ce{Mn^{2+}} de alto espín ($n = 5$) $5.92\,\mu_B$, \ce{Fe^{2+}} (4) $4.90\,\mu_B$, \ce{Co^{2+}} (3) $3.87\,\mu_B$; \ce{Fe^{2+}} y
\ce{Co^{3+}} de bajo espín ($t_{2g}^6$): 0, diamagnéticos.
\end{solution}

\begin{solution}{exo:b3:complex-spectra-magnetism:4}
$\ce{[Mn(H2O)6]^{2+}}$ (prohibida por el espín y por Laporte, casi incoloro) $<$ $\ce{[Co(H2O)6]^{2+}}$
(prohibida por Laporte, vibrónica) $<$ $\ce{[CoCl4]^{2-}}$ (sin centro de simetría) $<$ \ce{MnO4-} (transferencia
de carga permitida).
\end{solution}

\begin{solution}{exo:b3:complex-spectra-magnetism:5}
8500: $^3$A$_{2g}\to{}^3$T$_{2g}$, luego $\Delta_{\mathrm o} = \qty{8500}{cm^{-1}}$; 14100: $^3$T$_{1g}$(F); 24900: $^3$T$_{1g}$(P).
$15B = 39\,000 - 3 \times 8500$, $B = \qty{900}{cm^{-1}}$, $\beta = 900/1056 = 0.85$.
\end{solution}

\begin{solution}{exo:b3:complex-spectra-magnetism:6}
$\Delta_{\mathrm o} = \qty{17400}{cm^{-1}}$; resolviendo $7.5B + 1.5\Delta_{\mathrm o} - \frac12\sqrt{225B^2 - 18B\Delta_{\mathrm o} + \Delta_{\mathrm o}^2} = 24\,600$: $B = \qty{729}{cm^{-1}}$,
$\beta = 0.79$. $\nu_3 = \qty{38500}{cm^{-1}}$ (\qty{260}{nm}), en el ultravioleta.
\end{solution}

\begin{solution}{exo:b3:complex-spectra-magnetism:7}
Los electrones $d$ están más deslocalizados hacia los iones óxido de las gemas que hacia las
moléculas de agua: allí los enlaces Cr--O son más covalentes.
\end{solution}

\begin{solution}{exo:b3:complex-spectra-magnetism:8}
Fuerte ($e_g$ lleno de forma desigual): $d^9$, $d^4$ de alto espín, $d^7$ de bajo espín. Nula: $d^3$ ($t_{2g}^3$, término A) y
$d^6$ de bajo espín ($t_{2g}^6$). Débil: $d^6$ de alto espín ($t_{2g}^4e_g^2$, término T).
\end{solution}

\begin{solution}{exo:b3:complex-spectra-magnetism:9}
$d^7$ octaédrico de alto espín: término fundamental $^4$T$_{1g}$, un término T con \omterm{def:b3:complex-spectra-magnetism:moment}{contribución orbital},
y un momento muy superior al de solo espín, $3.87\,\mu_B$. $d^7$ tetraédrico ($\equiv$ $d^3$ octaédrico): término
fundamental $^4$A$_2$, momento orbital extinguido, próximo al de solo espín (algo superior, por
mezcla con términos T excitados).
\end{solution}

\begin{solution}{exo:b3:complex-spectra-magnetism:10}
$\mu_{\mathrm{eff}} = 2.828\sqrt{1.87} = 3.87\,\mu_B$: tres electrones desapareados ($S = \frac32$).
\end{solution}

\begin{solution}{exo:b3:complex-spectra-magnetism:11}
$\Delta\chi_v = 3 \times 25.0/\num{400e6} = \num{1.88e-7}$; $\chi_{\mathrm m} = \num{1.88e-7}/\qty{10.0}{mol.m^{-3}} = \qty{1.88e-8}{m^3.mol^{-1}}$, es decir,
\qty{1.49e-3}{cm^3.mol^{-1}}; $\chi_{\mathrm m}T = 0.445$, $\mu_{\mathrm{eff}} = 1.89\,\mu_B$: un electrón desapareado.
\end{solution}

\begin{solution}{exo:b3:complex-spectra-magnetism:12}
$T_{1/2} = 15\,000/75 = \qty{200}{K}$. A \qty{250}{K}: $\Delta G = 15\,000 - 250 \times 75 = \qty{-3750}{J/mol}$, $K = \eu^{1.80} = 6.07$,
$x_{\mathrm{HS}} = 0.86$; $\chi_{\mathrm m}T = 0.86 \times 3.00 = \qty{2.6}{cm^3.K.mol^{-1}}$.
\end{solution}

\begin{solution}{pb:b3:complex-spectra-magnetism:1}
\textbf{1.} [Ar]$3d^3$.
\textbf{2.} $^4$F.
\textbf{3.} $^4$A$_{2g}$ + $^4$T$_{2g}$ + $^4$T$_{1g}$.
\textbf{4.} $^4$A$_{2g}$, el término de $t_{2g}^3$, con los tres electrones en los orbitales inferiores.
\textbf{5.} $^4$T$_{1g}$(P).
\textbf{6.} $^4$A$_{2g}\to{}^4$T$_{2g}$, $\to{}^4$T$_{1g}$(F), $\to{}^4$T$_{1g}$(P).
\textbf{7.} 561 y \qty{414}{nm}.
\textbf{8.} 597 y \qty{434}{nm}.
\textbf{9.} El rubí absorbe el amarillo verdoso y el violeta y transmite el rojo (con algo
de azul); las bandas de la esmeralda, desplazadas hacia el rojo, absorben también el rojo anaranjado, y se
transmite el verde.
\textbf{10.} Rubí \qty{17820}{cm^{-1}}, esmeralda \qty{16740}{cm^{-1}}.
\textbf{11.} $^4$T$_{2g}$ ($t_{2g}^2e_g$) y $^4$A$_{2g}$ ($t_{2g}^3$) tienen la misma energía de repulsión electrónica; su
diferencia es la energía de promoción monoelectrónica $\Delta_{\mathrm o}$.
\textbf{12.} $B = (14\,305 - 547)/15 = \qty{917}{cm^{-1}}$.
\textbf{13.} \qty{614}{cm^{-1}}.
\textbf{14.} \qty{618}{cm^{-1}}.
\textbf{15.} 0.67 para ambas.
\textbf{16.} 29.0 y 27.1.
\textbf{17.} $\nu_3 = \qty{38500}{cm^{-1}}$ (\qty{260}{nm}) para el rubí, en el ultravioleta, donde la absorción de
transferencia de carga la oculta.
\textbf{18.} El sitio octaédrico del berilo es mayor: unos enlaces Cr--O más largos dan un
desdoblamiento menor, que varía muy deprisa con la distancia.
\textbf{19.} $10^7/695 = \qty{14390}{cm^{-1}}$, $23.4B$.
\textbf{20.} $^2$E$_g$ depende sobre todo de $C$: el nivel medido, más alto, indica que $C/B$ es mayor,
de unos 5.3 según la misma diagonalización.
\textbf{21.} $^2$E$_g$ y $^4$A$_{2g}$ proceden ambos de $t_{2g}^3$ y solo difieren por la inversión de un espín: el enlace,
y por tanto la geometría, es el mismo en ambos, y la transición no tiene progresión
vibracional (una raya 0--0).
\textbf{22.} Está prohibida por el espín.
\textbf{23.} Su energía depende de $B$ y de $C$, y casi nada de $\Delta_{\mathrm o}$ (una recta plana en el diagrama),
y $B$ es casi el mismo en las dos gemas.
\textbf{24.} Los iones bombeados a las bandas anchas se relajan y se acumulan en el nivel
$^2$E$_g$, de vida larga, de modo que puede haber más iones en él que en el nivel fundamental: una inversión de
\omterm{def:b3:partition-functions:boltzmann}{población}.
\textbf{25.} \textbf{$\Delta_{\mathrm o}(\text{rubí}) - \Delta_{\mathrm o}(\text{esmeralda}) = \qty{1080}{cm^{-1}}$}: toda la diferencia entre el rojo y el
verde.
\end{solution}
